% Options for packages loaded elsewhere
\PassOptionsToPackage{unicode}{hyperref}
\PassOptionsToPackage{hyphens}{url}
%
\documentclass[
]{article}
\usepackage{amsmath,amssymb}
\usepackage{lmodern}
\usepackage{iftex}
\ifPDFTeX
  \usepackage[T1]{fontenc}
  \usepackage[utf8]{inputenc}
  \usepackage{textcomp} % provide euro and other symbols
\else % if luatex or xetex
  \usepackage{unicode-math}
  \defaultfontfeatures{Scale=MatchLowercase}
  \defaultfontfeatures[\rmfamily]{Ligatures=TeX,Scale=1}
\fi
% Use upquote if available, for straight quotes in verbatim environments
\IfFileExists{upquote.sty}{\usepackage{upquote}}{}
\IfFileExists{microtype.sty}{% use microtype if available
  \usepackage[]{microtype}
  \UseMicrotypeSet[protrusion]{basicmath} % disable protrusion for tt fonts
}{}
\makeatletter
\@ifundefined{KOMAClassName}{% if non-KOMA class
  \IfFileExists{parskip.sty}{%
    \usepackage{parskip}
  }{% else
    \setlength{\parindent}{0pt}
    \setlength{\parskip}{6pt plus 2pt minus 1pt}}
}{% if KOMA class
  \KOMAoptions{parskip=half}}
\makeatother
\usepackage{xcolor}
\usepackage[margin=1in]{geometry}
\usepackage{color}
\usepackage{fancyvrb}
\newcommand{\VerbBar}{|}
\newcommand{\VERB}{\Verb[commandchars=\\\{\}]}
\DefineVerbatimEnvironment{Highlighting}{Verbatim}{commandchars=\\\{\}}
% Add ',fontsize=\small' for more characters per line
\usepackage{framed}
\definecolor{shadecolor}{RGB}{248,248,248}
\newenvironment{Shaded}{\begin{snugshade}}{\end{snugshade}}
\newcommand{\AlertTok}[1]{\textcolor[rgb]{0.94,0.16,0.16}{#1}}
\newcommand{\AnnotationTok}[1]{\textcolor[rgb]{0.56,0.35,0.01}{\textbf{\textit{#1}}}}
\newcommand{\AttributeTok}[1]{\textcolor[rgb]{0.77,0.63,0.00}{#1}}
\newcommand{\BaseNTok}[1]{\textcolor[rgb]{0.00,0.00,0.81}{#1}}
\newcommand{\BuiltInTok}[1]{#1}
\newcommand{\CharTok}[1]{\textcolor[rgb]{0.31,0.60,0.02}{#1}}
\newcommand{\CommentTok}[1]{\textcolor[rgb]{0.56,0.35,0.01}{\textit{#1}}}
\newcommand{\CommentVarTok}[1]{\textcolor[rgb]{0.56,0.35,0.01}{\textbf{\textit{#1}}}}
\newcommand{\ConstantTok}[1]{\textcolor[rgb]{0.00,0.00,0.00}{#1}}
\newcommand{\ControlFlowTok}[1]{\textcolor[rgb]{0.13,0.29,0.53}{\textbf{#1}}}
\newcommand{\DataTypeTok}[1]{\textcolor[rgb]{0.13,0.29,0.53}{#1}}
\newcommand{\DecValTok}[1]{\textcolor[rgb]{0.00,0.00,0.81}{#1}}
\newcommand{\DocumentationTok}[1]{\textcolor[rgb]{0.56,0.35,0.01}{\textbf{\textit{#1}}}}
\newcommand{\ErrorTok}[1]{\textcolor[rgb]{0.64,0.00,0.00}{\textbf{#1}}}
\newcommand{\ExtensionTok}[1]{#1}
\newcommand{\FloatTok}[1]{\textcolor[rgb]{0.00,0.00,0.81}{#1}}
\newcommand{\FunctionTok}[1]{\textcolor[rgb]{0.00,0.00,0.00}{#1}}
\newcommand{\ImportTok}[1]{#1}
\newcommand{\InformationTok}[1]{\textcolor[rgb]{0.56,0.35,0.01}{\textbf{\textit{#1}}}}
\newcommand{\KeywordTok}[1]{\textcolor[rgb]{0.13,0.29,0.53}{\textbf{#1}}}
\newcommand{\NormalTok}[1]{#1}
\newcommand{\OperatorTok}[1]{\textcolor[rgb]{0.81,0.36,0.00}{\textbf{#1}}}
\newcommand{\OtherTok}[1]{\textcolor[rgb]{0.56,0.35,0.01}{#1}}
\newcommand{\PreprocessorTok}[1]{\textcolor[rgb]{0.56,0.35,0.01}{\textit{#1}}}
\newcommand{\RegionMarkerTok}[1]{#1}
\newcommand{\SpecialCharTok}[1]{\textcolor[rgb]{0.00,0.00,0.00}{#1}}
\newcommand{\SpecialStringTok}[1]{\textcolor[rgb]{0.31,0.60,0.02}{#1}}
\newcommand{\StringTok}[1]{\textcolor[rgb]{0.31,0.60,0.02}{#1}}
\newcommand{\VariableTok}[1]{\textcolor[rgb]{0.00,0.00,0.00}{#1}}
\newcommand{\VerbatimStringTok}[1]{\textcolor[rgb]{0.31,0.60,0.02}{#1}}
\newcommand{\WarningTok}[1]{\textcolor[rgb]{0.56,0.35,0.01}{\textbf{\textit{#1}}}}
\usepackage{graphicx}
\makeatletter
\def\maxwidth{\ifdim\Gin@nat@width>\linewidth\linewidth\else\Gin@nat@width\fi}
\def\maxheight{\ifdim\Gin@nat@height>\textheight\textheight\else\Gin@nat@height\fi}
\makeatother
% Scale images if necessary, so that they will not overflow the page
% margins by default, and it is still possible to overwrite the defaults
% using explicit options in \includegraphics[width, height, ...]{}
\setkeys{Gin}{width=\maxwidth,height=\maxheight,keepaspectratio}
% Set default figure placement to htbp
\makeatletter
\def\fps@figure{htbp}
\makeatother
\setlength{\emergencystretch}{3em} % prevent overfull lines
\providecommand{\tightlist}{%
  \setlength{\itemsep}{0pt}\setlength{\parskip}{0pt}}
\setcounter{secnumdepth}{-\maxdimen} % remove section numbering
\ifLuaTeX
  \usepackage{selnolig}  % disable illegal ligatures
\fi
\IfFileExists{bookmark.sty}{\usepackage{bookmark}}{\usepackage{hyperref}}
\IfFileExists{xurl.sty}{\usepackage{xurl}}{} % add URL line breaks if available
\urlstyle{same} % disable monospaced font for URLs
\hypersetup{
  pdftitle={Untitled},
  pdfauthor={Vincent LE MEUR},
  hidelinks,
  pdfcreator={LaTeX via pandoc}}

\title{Untitled}
\author{Vincent LE MEUR}
\date{2023-09-07}

\begin{document}
\maketitle

\hypertarget{tp1-corruxe9lation-linuxe9aire-et-ruxe9gression-linuxe9aire-multiple-uxe0-laide-de-r}{%
\section{TP1 : Corrélation linéaire et Régression Linéaire Multiple à
l'aide de
R}\label{tp1-corruxe9lation-linuxe9aire-et-ruxe9gression-linuxe9aire-multiple-uxe0-laide-de-r}}

\hypertarget{compte-rendu}{%
\subsection{Compte rendu}\label{compte-rendu}}

\begin{Shaded}
\begin{Highlighting}[]
\FunctionTok{library}\NormalTok{(readxl)}
\NormalTok{TP1exel }\OtherTok{\textless{}{-}} \FunctionTok{read\_excel}\NormalTok{(}\StringTok{"TP1exel.xlsx"}\NormalTok{)}
\FunctionTok{View}\NormalTok{(TP1exel)}
\end{Highlighting}
\end{Shaded}

\hypertarget{etape-date}{%
\subsubsection{1 Etape Date}\label{etape-date}}

\begin{Shaded}
\begin{Highlighting}[]
\FunctionTok{head}\NormalTok{(TP1exel)}
\end{Highlighting}
\end{Shaded}

\begin{verbatim}
## # A tibble: 6 x 9
##   SOL      TR   A09   A12   A15   A18   A21   A24   A27
##   <chr> <dbl> <dbl> <dbl> <dbl> <dbl> <dbl> <dbl> <dbl>
## 1 SL01      2 0.632 0.292 0.318 0.436 0.296 0.069 0.079
## 2 SL02      4 0.552 0.275 0.418 0.468 0.258 0.116 0.072
## 3 SL03      6 0.565 0.3   0.392 0.501 0.279 0.04  0.052
## 4 SL04      8 0.549 0.332 0.502 0.509 0.224 0.055 0.018
## 5 SL05     10 0.57  0.351 0.449 0.48  0.222 0.056 0.025
## 6 SL06     12 0.273 0.309 0.427 0.324 0.156 0.056 0.08
\end{verbatim}

\begin{Shaded}
\begin{Highlighting}[]
\FunctionTok{summary}\NormalTok{(TP1exel)}
\end{Highlighting}
\end{Shaded}

\begin{verbatim}
##      SOL                  TR             A09              A12        
##  Length:16          Min.   : 2.00   Min.   :0.2540   Min.   :0.2750  
##  Class :character   1st Qu.: 9.50   1st Qu.:0.4677   1st Qu.:0.3202  
##  Mode  :character   Median :15.00   Median :0.5265   Median :0.4110  
##                     Mean   :15.38   Mean   :0.5098   Mean   :0.4689  
##                     3rd Qu.:22.50   3rd Qu.:0.5663   3rd Qu.:0.5835  
##                     Max.   :28.00   Max.   :0.7540   Max.   :0.7930  
##       A15              A18              A21              A24         
##  Min.   :0.3180   Min.   :0.2650   Min.   :0.0630   Min.   :0.01900  
##  1st Qu.:0.4195   1st Qu.:0.4510   1st Qu.:0.1482   1st Qu.:0.04975  
##  Median :0.5260   Median :0.5050   Median :0.1765   Median :0.05950  
##  Mean   :0.5638   Mean   :0.4956   Mean   :0.1949   Mean   :0.06463  
##  3rd Qu.:0.6730   3rd Qu.:0.5218   3rd Qu.:0.2587   3rd Qu.:0.07975  
##  Max.   :0.9390   Max.   :0.7850   Max.   :0.3130   Max.   :0.11600  
##       A27         
##  Min.   :0.00600  
##  1st Qu.:0.03475  
##  Median :0.07200  
##  Mean   :0.06031  
##  3rd Qu.:0.08075  
##  Max.   :0.10000
\end{verbatim}

\includegraphics{TP1_files/figure-latex/unnamed-chunk-4-1.pdf}

\hypertarget{calcul-de-la-matrice-des-corruxe9lations-linuxe9aires-et-analyse}{%
\subsubsection{Calcul de la Matrice des Corrélations linéaires et
analyse}\label{calcul-de-la-matrice-des-corruxe9lations-linuxe9aires-et-analyse}}

\begin{verbatim}
Le chargement a nécessité le package : lattice
\end{verbatim}

\begin{verbatim}
Le chargement a nécessité le package : survival
\end{verbatim}

\begin{verbatim}
Le chargement a nécessité le package : Formula
\end{verbatim}

\begin{verbatim}
Le chargement a nécessité le package : ggplot2
\end{verbatim}

\begin{verbatim}

Attachement du package : 'Hmisc'
\end{verbatim}

\begin{verbatim}
Les objets suivants sont masqués depuis 'package:base':

    format.pval, units
\end{verbatim}

\begin{verbatim}
       TR  A09  A12  A15  A18   A21  A24  A27
TR   1.00 0.23 0.96 0.92 0.59 -0.23 0.02 0.40
A09  0.23 1.00 0.48 0.56 0.88  0.83 0.25 0.37
A12  0.96 0.48 1.00 0.97 0.76  0.02 0.08 0.51
A15  0.92 0.56 0.97 1.00 0.83  0.13 0.11 0.47
A18  0.59 0.88 0.76 0.83 1.00  0.61 0.22 0.46
A21 -0.23 0.83 0.02 0.13 0.61  1.00 0.21 0.27
A24  0.02 0.25 0.08 0.11 0.22  0.21 1.00 0.30
A27  0.40 0.37 0.51 0.47 0.46  0.27 0.30 1.00

n= 16 


P
    TR     A09    A12    A15    A18    A21    A24    A27   
TR         0.3946 0.0000 0.0000 0.0164 0.3967 0.9556 0.1285
A09 0.3946        0.0617 0.0251 0.0000 0.0000 0.3598 0.1527
A12 0.0000 0.0617        0.0000 0.0006 0.9424 0.7605 0.0437
A15 0.0000 0.0251 0.0000        0.0000 0.6187 0.6916 0.0636
A18 0.0164 0.0000 0.0006 0.0000        0.0113 0.4129 0.0760
A21 0.3967 0.0000 0.9424 0.6187 0.0113        0.4409 0.3097
A24 0.9556 0.3598 0.7605 0.6916 0.4129 0.4409        0.2616
A27 0.1285 0.1527 0.0437 0.0636 0.0760 0.3097 0.2616       
\end{verbatim}

\begin{verbatim}
corrplot 0.92 loaded
\end{verbatim}

\includegraphics{TP1_files/figure-latex/unnamed-chunk-7-1.pdf}

On regarde les correlations dont la p valeur est différente de 0

\end{document}
